\documentclass[10pt]{article}
\usepackage{amsmath}
\usepackage{amssymb}
\usepackage{hyperref}

\usepackage{tikz}
\usetikzlibrary{positioning, arrows.meta}


\title{E.A.6.4 - nQueens SAT}
\author{}
\date{}

\begin{document}
\maketitle
\vspace{-6em}

\section{Modellazione}

Sia $x_{i,j}$ una variabile proposizionale tale che $x_{i,j} = \top \iff$ c'è una regina nella posizione $[i, j]$ della scacchiera $n \times n$. 

Introduciamo le seguenti formule:

\paragraph{Almeno una regina per riga:}
\[
    \varphi_{ar} = \bigwedge_{i=1}^{n} \bigvee_{j=1}^{n} x_{i,j}
\]

\paragraph{Al massimo una regina per riga:}
\[
\varphi_{mr} =\bigwedge_{i=1}^{n} \bigwedge_{j_1=1}^{n-1} \bigwedge_{j_2=j_1+1}^{n} (\neg x_{i,j_1} \lor \neg x_{i,j_2})
\]

\paragraph{Al massimo una regina per colonna:}
\[
    \varphi_{mc}=\bigwedge_{j=1}^{n} \bigwedge_{i_1=1}^{n-1} \bigwedge_{i_2=i_1+1}^{n} (\neg x_{i_1,j} \lor \neg x_{i_2,j})
\]

\paragraph{Al massimo una regina per diagonale discendente:}
\[
    \varphi_{md1}=\bigwedge_{i=1}^{n-1} \bigwedge_{j=1}^{n-1} \bigwedge_{k=1}^{\min(n-i, n-j)} (\neg x_{i,j} \lor \neg x_{i+k, j+k})
\]

\paragraph{Al massimo una regina per diagonale ascendente:}
\[
\varphi_{md2}=\bigwedge_{i=1}^{n-1} \bigwedge_{j=2}^{n} \bigwedge_{k=1}^{\min(n-i, j-1)} (\neg x_{i,j} \lor \neg x_{i+k, j-k})
\]

\paragraph{}

La codifica in CNF del problema n-queens è data dalla loro congiunzione:
$$\Phi = \varphi_{ar} \land \varphi_{mc} \land \varphi_{mr} \land \varphi_{md1} \land \varphi_{md2}$$

\section{Istanziazione}
Le variabili sono $x_{i,j}$ con $i, j \in \{1, 2, 3, 4\}$. $\Phi$ è la congiunzione dei seguenti insiemi di clausole.

\paragraph{1. Almeno una regina per riga (4 clausole):}
\[
\begin{aligned}
    & (x_{1,1} \lor x_{1,2} \lor x_{1,3} \lor x_{1,4}) \land \\
    & (x_{2,1} \lor x_{2,2} \lor x_{2,3} \lor x_{2,4}) \land \\
    & (x_{3,1} \lor x_{3,2} \lor x_{3,3} \lor x_{3,4}) \land \\
    & (x_{4,1} \lor x_{4,2} \lor x_{4,3} \lor x_{4,4})
\end{aligned}
\]

\paragraph{2. Al massimo una regina per riga (24 clausole, 6 per riga):}
Per la riga $i=1$:
\[
\begin{aligned}
    & (\neg x_{1,1} \lor \neg x_{1,2}) \land (\neg x_{1,1} \lor \neg x_{1,3}) \land (\neg x_{1,1} \lor \neg x_{1,4}) \land \\
    & (\neg x_{1,2} \lor \neg x_{1,3}) \land (\neg x_{1,2} \lor \neg x_{1,4}) \land (\neg x_{1,3} \lor \neg x_{1,4})
\end{aligned}
\]
\textit{(queste 6 combinazioni si ripetono per le righe $2, 3, 4$)}

\paragraph{3. Al massimo una regina per colonna (24 clausole, 6 per colonna):}
Per la colonna $j=1$:
\[
\begin{aligned}
    & (\neg x_{1,1} \lor \neg x_{2,1}) \land (\neg x_{1,1} \lor \neg x_{3,1}) \land (\neg x_{1,1} \lor \neg x_{4,1}) \land \\
    & (\neg x_{2,1} \lor \neg x_{3,1}) \land (\neg x_{2,1} \lor \neg x_{4,1}) \land (\neg x_{3,1} \lor \neg x_{4,1})
\end{aligned}
\]
\textit{(si ripetono per le colonne $2, 3, 4$)}

\paragraph{4. Al massimo una regina per diagonale discendente $\searrow$ (14 clausole):}
\begin{itemize}
    \item Diagonale principale (4 celle: $x_{1,1}, x_{2,2}, x_{3,3}, x_{4,4}$): genera 6 clausole.
    \[
    \begin{aligned}
        & (\neg x_{1,1} \lor \neg x_{2,2}) \land (\neg x_{1,1} \lor \neg x_{3,3}) \land (\neg x_{1,1} \lor \neg x_{4,4}) \land \\
        & (\neg x_{2,2} \lor \neg x_{3,3}) \land (\neg x_{2,2} \lor \neg x_{4,4}) \land (\neg x_{3,3} \lor \neg x_{4,4})
    \end{aligned}
    \]
    \item Sopra la d. principale (3 celle: $x_{1,2}, x_{2,3}, x_{3,4}$): genera 3 clausole.
    \[
        (\neg x_{1,2} \lor \neg x_{2,3}) \land (\neg x_{1,2} \lor \neg x_{3,4}) \land (\neg x_{2,3} \lor \neg x_{3,4})
    \]
    \item Sotto la d. principale (3 celle: $x_{2,1}, x_{3,2}, x_{4,3}$): genera 3 clausole.
    \[
        (\neg x_{2,1} \lor \neg x_{3,2}) \land (\neg x_{2,1} \lor \neg x_{4,3}) \land (\neg x_{3,2} \lor \neg x_{4,3})
    \]
    \item D. agli estremi (2 celle ciascuna): generano 1 clausola ciascuna.
    \[
        (\neg x_{1,3} \lor \neg x_{2,4}) \quad \text{e} \quad (\neg x_{3,1} \lor \neg x_{4,2})
    \]
\end{itemize}

\paragraph{5. Al massimo una regina per diagonale ascendente $\nearrow$ (14 clausole):}
\begin{itemize}
    \item Diagonale principale (4 celle: $x_{4,1}, x_{3,2}, x_{2,3}, x_{1,4}$): genera 6 clausole.
    \[
    \begin{aligned}
        & (\neg x_{4,1} \lor \neg x_{3,2}) \land (\neg x_{4,1} \lor \neg x_{2,3}) \land (\neg x_{4,1} \lor \neg x_{1,4}) \land \\
        & (\neg x_{3,2} \lor \neg x_{2,3}) \land (\neg x_{3,2} \lor \neg x_{1,4}) \land (\neg x_{2,3} \lor \neg x_{1,4})
    \end{aligned}
    \]
    \item Sopra la d. principale (3 celle: $x_{3,1}, x_{2,2}, x_{1,3}$): genera 3 clausole.
    \[
        (\neg x_{3,1} \lor \neg x_{2,2}) \land (\neg x_{3,1} \lor \neg x_{1,3}) \land (\neg x_{2,2} \lor \neg x_{1,3})
    \]
    \item Sotto la d. principale (3 celle: $x_{4,2}, x_{3,3}, x_{2,4}$): genera 3 clausole.
    \[
        (\neg x_{4,2} \lor \neg x_{3,3}) \land (\neg x_{4,2} \lor \neg x_{2,4}) \land (\neg x_{3,3} \lor \neg x_{2,4})
    \]
    \item D. agli estremi (2 celle ciascuna): generano 1 clausola ciascuna.
    \[
        (\neg x_{2,1} \lor \neg x_{1,2}) \quad \text{e} \quad (\neg x_{4,3} \lor \neg x_{3,4})
    \]
\end{itemize}

\section{Codec}
Ho scritto un piccolo encoder molto naive in Haskell [\href{https://github.com/AglaiaNorza/ia/blob/main/SATencoder/naive-encoder.hs}{qui}] e uno meno naive (multithreaded) in c++ [\href{https://github.com/AglaiaNorza/ia/tree/main/SATencoder/multithreaded}{qui}] (ho inserito anche encoding tramite Tseitin transformation per formule non nativamente in CNF).

\end{document}
